\subsection{Schnorr knowledge proofs and explicit OR composition}\label{sec:realize:or}
The three-move proof of knowledge for discrete logarithms
is due to Schnorr~\cite{Schnorr1991}; its OR composition is
a standard instance of the partial-knowledge methods
of Cramer and coauthors~\cite{CramerDamgardSchoenmakers1994}.
We give all protocol messages and their elementary proofs.

\begin{definition}[Schnorr relation]\label{def:realize:schnorr}
For a group element $Y\in G$, define the relation
$\mathcal R_{\log}(Y;r)$ by $Y=h^r$,
$r\in\mathbb Z_P$.
The prover sends $T=h^u$ for uniform
$u\in\mathbb Z_P$; the verifier responds with
uniform challenge $e\in\mathbb Z_P$; the
prover returns $z=u+er$.
The verifier accepts if
$h^z=TY^e$.
\end{definition}

\begin{lemma}[Schnorr completeness]\label{lem:realize:schnorr-complete}
For every valid witness $r$ and all $u,e$,
the Schnorr verifier accepts.
\end{lemma}
\begin{proof}
Since $Y=h^r$ and $z=u+er$,
$h^z=h^uh^{er}=TY^e$.
\end{proof}

\begin{lemma}[Two Schnorr transcripts extract the logarithm]\label{lem:realize:schnorr-extract}
If two accepting Schnorr transcripts have the same
$Y,T$ but distinct challenges $e\ne e'$ and
responses $z,z'$, then
\[
 \boxed{r=(z-z')(e-e')^{-1}}
\]
satisfies $Y=h^r$.
\end{lemma}
\begin{proof}
Divide the two verifier equations to obtain
$h^{z-z'}=Y^{e-e'}$.
Since $e-e'\ne0$ modulo the prime $P$, it has
an inverse. Raise both sides to that inverse:
$h^{(z-z')(e-e')^{-1}}=Y$.
\end{proof}

\begin{lemma}[Perfect honest-verifier Schnorr simulation]\label{lem:realize:schnorr-hvzk}
For a uniformly sampled verifier challenge $e$,
the joint transcript has exactly the distribution obtained
without the witness by sampling independent uniform
$(e,z)$ and setting $T=h^zY^{-e}$.
\end{lemma}
\begin{proof}
For any fixed valid witness $r$,
the honest challenge $e$ is uniform and independent of
uniform prover nonce $u$. Conditional on $e$,
the response $z=u+er$ is uniform on $\mathbb Z_P$.
The verifier equation uniquely determines
$T=h^zY^{-e}$ from $(e,z)$.
Therefore the simulated and honest distributions
of the entire triple coincide.
\end{proof}

\begin{definition}[Schnorr OR proof for a finite candidate set]\label{def:realize:or}
Fix nonempty $S=\{t_1,\ldots,t_d\}\subseteq\mathbb Z_P$
with distinct elements and a public commitment $C\in G$.
For each branch set $Y_j=Cg^{-t_j}$.
The prover knows an index $j_*$ and a scalar $r$
such that $Y_{j_*}=h^r$.
Before receiving the global verifier challenge,
it chooses independent uniform
$e_j,z_j\in\mathbb Z_P$ for $j\ne j_*$,
sets
\[
 T_j=h^{z_j}Y_j^{-e_j}\quad(j\ne j_*),
\]
and chooses uniform $u$ to set $T_{j_*}=h^u$.
It sends $(T_1,\ldots,T_d)$.
After the verifier draws uniform $e\in\mathbb Z_P$,
the prover sets
\[
 e_{j_*}:=e-\sum_{j\ne j_*}e_j,\qquad
 z_{j_*}:=u+e_{j_*}r,
\]
and sends all $(e_j,z_j)_{j=1}^d$.
The verifier accepts if
\[
 \sum_j e_j=e
 \quad\text{and}\quad
 h^{z_j}=T_jY_j^{e_j}\ \ \forall j.
\]
All challenge and response scalars belong to $\mathbb Z_P$.
\end{definition}

\begin{lemma}[OR-proof perfect completeness]\label{lem:realize:or-complete}
For any candidate $t_{j_*}$ with a known
opening $C=g^{t_{j_*}}h^r$, every honest
execution of Definition~\ref{def:realize:or} accepts.
\end{lemma}
\begin{proof}
By construction the local challenges sum to the
verifier's global challenge.
For $j\ne j_*$ the prover defined $T_j$ precisely so that
$h^{z_j}=T_jY_j^{e_j}$.
For $j=j_*$,
$T_{j_*}Y_{j_*}^{e_{j_*}}
=h^uh^{re_{j_*}}=h^{u+e_{j_*}r}=h^{z_{j_*}}$.
Thus all equations pass.
\end{proof}

\begin{theorem}[OR-proof perfect honest-verifier zero knowledge]\label{thm:realize:or-hvzk}
With uniform independent global challenge $e$,
the full accepting OR-proof transcript can be
simulated \emph{exactly} without any opening of $C$:
choose independent uniform $(e_j,z_j)\in\mathbb Z_P^2$
for every branch, put $e=\sum_j e_j$, and set
\[
 \boxed{T_j=h^{z_j}(Cg^{-t_j})^{-e_j}.}
\]
The resulting joint distribution agrees with
the honest prover/verifier distribution for
every valid witness branch.
\end{theorem}
\begin{proof}
For the honest prover, the challenges $e_j$ on
false branches are independent uniform.
The true-branch challenge
$e_{j_*}=e-\sum_{j\ne j_*}e_j$ is uniform
and independent of them because the global
honest-verifier challenge is uniform.
The false-branch responses $z_j$ were drawn
independent uniform; the true response
$z_{j_*}=u+e_{j_*}r$ is also uniform
conditioned on the entire challenge vector
because $u$ was uniform independent.
Thus the complete vector $(e_j,z_j)_j$
is uniform on $(\mathbb Z_P^2)^d$.
Every first message $T_j$ is uniquely
determined by the verifier equation as
$h^{z_j}Y_j^{-e_j}$.
The simulator uses exactly this deterministic
expression from exactly the same distribution,
so the complete transcript laws coincide.
\end{proof}

\begin{theorem}[OR-proof special soundness]\label{thm:realize:or-extract}
Two accepting OR-proof transcripts for
the same statement $C,S$ and the same first-message
vector $(T_j)_j$ but different global
challenges $e\ne e'$ determine an index $j$
and an opening
\[
 \boxed{C=g^{t_j}h^r,\qquad
 r=(z_j-z_j')(e_j-e_j')^{-1}.}
\]
\end{theorem}
\begin{proof}
The sums of the two local challenge vectors
are different, so for some branch $j$
we have $e_j\ne e_j'$.
Divide the two accepting verifier equations
for that branch, obtaining
$h^{z_j-z_j'}=Y_j^{e_j-e_j'}$.
As in Lemma~\ref{lem:realize:schnorr-extract},
inverting the nonzero challenge difference yields
$Y_j=h^r$ for the displayed $r$.
Since $Y_j=Cg^{-t_j}$, multiplying by $g^{t_j}$
gives $C=g^{t_j}h^r$.
\end{proof}

\begin{remark}[Proof of knowledge versus a nontrivial language]
Because $h$ generates the entire group $G$,
each group element $Y_j$ has some logarithm
mathematically. The OR protocol proves
\emph{efficient knowledge of an opening},
not mere set membership. Its computational
knowledge-soundness guarantee requires
the standard rewinding extractor and
independent uniform global challenges;
the two-transcript special-soundness
theorem does not alone bound a malicious
prover's success probability in the
entire composed protocol.
\end{remark}

\begin{remark}[Honest-verifier scope]
Theorem~\ref{thm:realize:or-hvzk} assumes the
sigma verifier chooses its challenge uniformly.
It makes no assertion of straight-line
zero knowledge against an arbitrary
malicious verifier that chooses the
sigma challenge as a function of its first
message. A separate compiler is needed for
that stronger notion.
\end{remark}
